% ==============================================================================
% Instituto Federal de Educa��o, Ci�ncia e Tecnologia do Rio Grande do Norte
% Modelo usado para Trabalho
%
% Autor:  Rodrigo Siqueira
% Baseado no modelo do PPGEE da UFRN feito pelo Prof. Adelardo
% ==============================================================================

% Classe .......................................................................
\documentclass[a4paper,12pt,openany,oneside]{book}


% Pacotes ....................................................................................
% Pacotes Principais -----------------------------------------------------------
\usepackage[portuges,brazil]{babel}
\usepackage{ae} 
%%\usepackage[utf8]{inputenc}
%\usepackage[latin1]{inputenc}

%NO MAC ao inv�s de latin1 deve-se usar o applemac
\usepackage[applemac]{inputenc}


% Formatacaoo de capitulos ------------------------------------------------------
%\usepackage{capitulos}

% Figuras e Imagens ------------------------------------------------------------
\usepackage{graphicx}
% Figuras lado a lado
%\usepackage{epsfig} - REMOVIDO
%\usepackage{subfigure} - REMOVIDO

% Para poder ter tabelas com colunas de largura auto-ajust�vel
\usepackage{tabularx}

% Utilizar H para inserir as imagens REALMENTE onde eu desejo
%\usepackage{float} - REMOVIDO

% Para que os primeiros par�grafos das se��es tamb�m sejam indentados
% \usepackage{indentfirst} - REMOVIDO

%Permitir usar o comando \citeasnoun que coloca o nome na cita��o
\usepackage{harvard}


% Fontes -----------------------------------------------------------------------
%\usepackage[T1]{fontenc} - REMOVIDO
%\usepackage{pslatex} - REMOVIDO

% Matem�tico -------------------------------------------------------------------
%\usepackage{amsmath} - REMOVIDO
%\usepackage{amstext} - REMOVIDO

% Simbolos ---------------------------------------------------------------------
%\usepackage{textcomp} - REMOVIDO

% Tabelas ----------------------------------------------------------------------
%\usepackage{multirow} - REMOVIDO
% Colorir a tabela
%\usepackage{colortbl} - REMOVIDO

% Gloss�rio --------------------------------------------------------------------
%\usepackage[portuguese,noprefix]{nomencl}
%\usepackage{makeglo}

% Outros pacotes ---------------------------------------------------------------
%\usepackage{noitemsep} - REMOVIDO

%Subsubsection
%\usepackage{chngcntr} - REMOVIDO

% Para permitir espa�amento simples, 1 1/2 e duplo
\usepackage{setspace}

% Para executar um comando depois do fim da p�gina corrente
%\usepackage{afterpage} - REMOVIDO

% Para poder incluir arquivos Postscript com cores (do Xfig, por exemplo)
%\usepackage{color} - REMOVIDO

  % Itens numerados
%\usepackage{enumerate} - REMOVIDO

% Coment�rios em bloco
%\usepackage{verbatim} - REMOVIDO

% Refer�ncias ------------------------------------------------------------------
%\usepackage{html} - REMOVIDO
%\usepackage{url} - REMOVIDO
%\usepackage[abbr]{harvard}	% As chamadas s�o sempre abreviadas - REMOVIDO
%\harvardparenthesis{square}	% Colchetes nas chamadas - REMOVIDO
%\harvardyearparenthesis{round}	% Par�ntesis nos anos das refer�ncias - REMOVIDO
%\renewcommand{\harvardand}{e}	% Substituir "&" por "e" nas refer�ncias
%\renewcommand{\harvardurl}{URL: \url} - REMOVIDO

% Novos Comandos .....................................................................
% newcommand define novos comandos, que podem passar a ser usados da
% mesma forma que os comandos LaTeX de base.

% Configura��o da fonte
%\renewcommand{\familydefault}{\sfdefault}

% Margens ----------------------------------------------------------------------
%\setlength{\oddsidemargin}{3.5cm}
%\setlength{\evensidemargin}{2.5cm}
%\setlength{\textwidth}{15cm}
%\addtolength{\oddsidemargin}{-1in}
%\addtolength{\evensidemargin}{-1in}
%%
%\setlength{\topmargin}{2.0cm}
%\setlength{\headheight}{1.0cm}
%\setlength{\headsep}{1.0cm}
%\setlength{\textheight}{22.7cm}
%\setlength{\footskip}{1.0cm}
%\addtolength{\topmargin}{-1in}

% Cap�tulos --------------------------------------------------------------------
% N�o aparecer o n�mero na primeira p�gina dos cap�tulos
\newcommand{\mychapter}[1]{\chapter{#1}\thispagestyle{empty}}
\newcommand{\mychapterstar}[1]{\chapter*{#1}\thispagestyle{empty}}

% Os cap�tulos sem n�mero
%\newcommand{\mychapterast}[1]{\chapter*{#1}\thispagestyle{empty}
%\chaptermark{#1}
%\afterpage{\markboth{\uppercase{#1}}{\rightmark}}
%\markboth{\uppercase{#1}}{}
%}

%Referencias
%\newcommand{\BibTeX}{\textsc{B\hspace{-0.1em}i\hspace{-0.1em}b\hspace{-0.3em}}\TeX}

% Comandos matem�ticos ---------------------------------------------------------
% Implica��o em f�rmulas
%\newcommand{\implica}{\quad\Rightarrow\quad} %Meio de linha
%\newcommand{\implicafim}{\quad\Rightarrow}   %Fim de linha
%\newcommand{\tende}{\rightarrow}

% Fra��o com parenteses
%\newcommand{\pfrac}[2]{\parent{\frac{#1}{#2}}}

% Transformada de Laplace e transformada Z
%\newcommand{\lapl}{\pounds}
%\newcommand{\transfz}{\mathcal{Z}}

% Sequ�ncias
%\newcommand{\sequencia}[4]{$#1_{#2}$, $#1_{#3}$, \ldots, $#1_{#4}$}

% Se��es sem n�mero
%\newcommand{\mysectionast}[1]{\section*{#1}
%\addcontentsline{toc}{section}{#1}
%\markright{\uppercase{#1}}
%}

% No tabularx, as celulas devem ser centradas verticalmente
\renewcommand{\tabularxcolumn}[1]{m{#1}}

% C�lulas centralizadas horizontalmente no tabularx
\newcolumntype{C}{>{\centering\arraybackslash}X}


% Outros ----------------------------------------------------------------------
%\newcommand{\chave}[1]{\left\{#1\right\}}
%\newcommand{\colchete}[1]{\left[#1\right]}
%\newcommand{\parent}[1]{\left(#1\right)}

%%subsubsection
%\newcounter{subcount}
%\newenvironment{mysub}
%{\begin{list}
%{\textbf{\thesubsubsection.\arabic{subcount}}}
%{\setlength{\itemindent}{2.2em}
%\setlength{\rightmargin}{.6in}
%\setlength{\labelwidth}{1in}
%\setlength{\labelsep}{.2in}
%\setlength{\parsep}{.5ex plus .2ex minus .1ex}
%\setlength{\itemsep}{0ex plus .2ex minus 0ex}
%\usecounter{subcount}}
%}


%\citationmode{abbr}

% Inicia o texto
\begin{document}

% Paginas iniciais (sem numera��o)
\pagestyle{empty}

% P�gina de rosto (capa interna)
%
% ********** P�gina de Rosto
%

% titlepage gera p�ginas sem numera��o
\begin{titlepage}

\begin{center}

\small

% O comando @{} no ambiente tabular x � para criar um novo delimitador
% entre colunas que n�o a barra vertical | que � normalmente utilizada.
% O delimitador desejado vai entre as chaves. No exemplo, n�o h� nada,
% de modo que o delimitador � vazio. Este recurso est� sendo usado para
% eliminar o espa�o que geralmente existe entre as colunas
\begin{tabularx}{\linewidth}{ c X }
% A figura foi colocada dentro de um parbox para que fique verticalmente
% centralizada em rela��o ao resto da linha
\parbox[c]{3cm}{\includegraphics[width=\linewidth]{IFRN}} &
\begin{center}
\textsf{\textsc{Instituto Federal de Educa��o Ci�ncia\\
 e Tecnologia do Rio Grande do Norte \\
Diretoria de Gest�o de Tecnologia da Informa��o \\
Curso Superior de Tecnologia em Redes de Computadores
}} 
\end{center}

\end{tabularx}


% O vfill � um espa�o vertical que assume a m�xima dimens�o poss�vel
% Os vfill's desta p�gina foram utilizados para que o texto ocupe
% toda a folha
\vfill

\LARGE

\textbf{Aqui � o t�tulo de seu trabalho}

\vfill

\Large

\textbf{Fulano do Anz�is Pereira}

\vfill

\normalsize

Prof. Dr. Rodrigo Siqueira Martins

\vfill

\hfill
\parbox{0.5\linewidth}{\textbf{%
% Descomente as op��es que se aplicam ao seu caso
%Proposta de Tema para Qualifica��o}
%Disserta��o de Mestrado}
Relat�rio} apresentado � Disciplina de Administra��o de Sistemas Propriet�rios
como parte dos requisitos para obten��o da nota da 1a unidade}


\vfill

\large

%Data
Natal, RN, junho de 2014

\end{center}

\end{titlepage}
  

%Espa�amento 1 1/2
\onehalfspacing  

% Paginas introdut�rias (com numera��o romana)
\frontmatter

% Lista de conte�do (sum�rio, gerado automaticamente)
\addcontentsline{toc}{chapter}{Sum�rio}   
\tableofcontents  

% Lista de figuras (gerada automaticamente)
\cleardoublepage 
\addcontentsline{toc}{chapter}{Lista de Figuras}  
\listoffigures  

% Lista de tabelas (gerada automaticamente)
\cleardoublepage 
\addcontentsline{toc}{chapter}{Lista de Tabelas}   
\listoftables  

% P�ginas do texto principal (com cabe�alho)
\mainmatter
\pagestyle{headings}

% Para facilitar a organiza��o, foi criado um diret�rio para cada
% cap�tulo do documento, pois assim os arquivos das figuras ficam
% classificados por cap�tulos

% Introdu��o
%%
%% Cap�tulo 1: Modelo de Cap�tulo
%%

% Est� sendo usando o comando \mychapter, que foi definido no arquivo
% comandos.tex. Este comando \mychapter � essencialmente o mesmo que o
% comando \chapter, com a diferen�a que acrescenta um \thispagestyle{empty}
% ap�s o \chapter. Isto � necess�rio para corrigir um erro de LaTeX, que
% coloca um n�mero de p�gina no rodap� de todas as p�ginas iniciais dos
% cap�tulos, mesmo quando o estilo de numera��o escolhido � outro.
\mychapter{Introdu��o}
\label{Cap:introducao}

%O que escrever
A introdu��o dever� ser escrita depois do desenvolvimento do trabalho. Nela deve conter a explica��o do que ser� abordado no trabalho e se necess�rio um hist�rico da necessidade do conte�do. Aqui � importante colocar cita��o direta onde se pode dizer que Segundo \citeasnoun{Bernardete} e ainda escrever seu pr�prio texto atrav�s da leitura de outros trabalhos e fazer uma cita��o indireta \cite{Bernardete}. N�o � muito interessante colocar figuras aqui na introdu��o embora n�o seja proibido. Procure escrever de forma resumida, mas fa�a valer ao leitor o que de importante ele encontrar� na leitura, ou seja, aqui deve-se "vender o peixe" do trabalho.

%Sobre esse trabalho
Descreva o m�ximo do que se trata o trabalho e sua import�ncia.

\section{Motiva��o}
\label{Sec:Motiva��o}
%Motiva��o
Explique aqui porque iniciou as pesquisas no tema e qual a motiva��o de desenvolver esse trabalho.

\section{Objetivo}
\label{Sec:Objetivo}
%Objetivo
Descreva o objetivo principal do trabalho. Tente criar um par�grafo resumindo o que seu trabalho far�. Depois seja mais espec�fico, pode inclusive criar t�picos para a Se��o \ref{Sec:ObjetivoEspecifico}.

\subsection{Objetivo espec�fico}
\label{Sec:ObjetivoEspecifico}
\begin{itemize}
  \item Criar uma ferramenta que ... ;
  \item Criar ... ;
\end{itemize}



\section{Organiza��o do trabalho}
\label{Sec:Organiza��o}
%Organiza��o do trabalho   
Sempre que precisar referenciar outra parte de seu trabalho use o comando \\ref apontando para o que voc� colocou no \\label como por exemplo essa Se��o \ref{Sec:Organiza��o}. Aqui voc� dever� informar como e porque seu trabalho foi organizado, informando o que ser� abordado em cada cap�tulo.


 

% Desenvolvimento
%%
%% Cap�tulo 2: Desenvolvimento do trabalho
%%

\mychapter{Desenvolvimento}
\label{Cap:desenvolvimento}

Este cap�tulo apresenta considera��es de ordem geral sobre a
organiza��o e desenvolvimento do texto.

\section{O que escrever}
\label{Sec:oque}

� poss�vel iniciar o cap�tulo de desenvolvimento informando como ser� desenvolvido o trabalho e o detalhe das ferramentas ou o conte�do que ser� abordado e utilizado. Em seguida comente o conte�do em conjunto com o conhecimento adquirido junto as consultas e refer�ncias bibliogr�ficas. 
Isso demonstra que foi feito uma pesquisa anterior e mostra que o trabalho tem relev�ncia cient�fica, ou seja, n�o est� sendo "inventada a roda". 

\section{Divis�es do documento e refer�ncias cruzadas}
\label{Sec:divisoes}

Sempre que for modificar/aprofundar partes de um tema, � interessante criar uma se��o de maneira a tornar o assunto mais focado e facilitar a leitura. 
N�o crie se��es sem que estas tenham nenhum texto ou sigam direto para outra se��o. Se isso acontecer � melhor rever como est� organizado o texto pois 
pode se tornar necess�rio remover cap�tulos ou se��es. 

No final de cada cap�tulo � importante apresentar coment�rio do que foi desenvolvido no texto permitindo ao leitor compreender a vis�o geral do conte�do e pode-se inclusive contextualizar o cap�tulo seguinte.

%%%%PAREI AQUI

Os cap�tulo, se��es e subse��es devem ser rotulados para que possam ser
referenciados em qualquer parte do texto.  Isto � feito com o comando
\verb|\label{}|, que deve ser colocado logo ap�s (nunca antes) o
comando que criou a se��o, cap�tulo, etc. O par�metro do comando
\texttt{label} � o nome simb�lico que ser� usado para se fazer
refer�ncia a esta entidade dentro do texto, com o comando
\verb|\ref{}|. O nome pode ser qualquer coisa, mas n�o pode conter
acentos, por exemplo. 

Veja a seguir um exemplo de
utiliza��o das refer�ncias cruzadas:
quotation � um ambiente para cita��es, que ficam "recuadas" em
rela��o ao resto do texto
\begin{quotation}
\dots no cap�tulo \ref{Cap:introducao} apresentamos um modelo de
cap�tulo do texto.
\end{quotation}

\section{Se��es}
\label{Sec:secoes}  %% labels n�o devem conter caracteres acentuados

S�o as divis�es dos cap�tulos. Essas se��es s�o do criadas vindo do geral para o mais espec�fico. Por exemplo de sistemas operacionais para processos, compartilhamento de mem�ria, etc.

\subsection{Subse��es}
\label{Sec:subsecoes}

Subse��es s�o divis�es das se��es. No exemplo do texto sobre
Processos pode-se criar uma subse��o de concorr�ncia, deadlock, etc

\subsubsection*{Subsubse��es}
\label{Sec:subsubsecoes}

Subsubse��es s�o divis�es de subse��es e n�o devem ser numeradas no
texto. 

\section{�ndices}

O \LaTeX\ � capaz de gerar automaticamente o �ndice do texto
(sum�rio), os indices de figuras e de tabelas.

\subsection{Sum�rio}
\label{Sec:sumario}

Todas as divis�es numeradas (cap�tulos, se��es e subse��es) s�o
automaticamente inclu�das no sum�rio. 

\subsection{Listas de figuras e tabelas}
\label{Sec:listasfigtab}

Estas listas s�o geradas automaticamente a partir dos
\texttt{caption}'s de todos os ambientes \texttt{figure} e
\texttt{table}. 

\section{Bibliografia}
\label{Sec:citacoes}

As refer�ncias bibliogr�ficas s�o inclu�das dentro de um ambiente
espec�fico para este fim, atrav�s dos comandos
\verb|\begin{thebibliography}| e \verb|\end{thebibliography}|. Um
documento s� pode conter um �nico destes ambientes. A refer�ncia aos
documentos no texto � feita usando-se refer�ncias cruzadas e chaves
simb�licas, da mesma forma que as equa��es e as figuras. Cada
documento dentro de um ambiente \texttt{thebibliography} � introduzido
por um comando \verb|\bibitem|. O argumento obrigat�rio (entre chaves)
do comando \texttt{bibitem} � a chave simb�lica pela qual o documento
ser� citado no texto, usando o comando \verb|\cite|.  

A principal vantagem da gera��o autom�tica de bibliografias � que se
pode manter um arquivo �nico com todas as refer�ncias bibliogr�ficas
que foram ou podem vir a ser usadas em algum dos seus documentos

Uma recomenda��o importante � que as cita��es n�o fazem parte do
texto; portanto, as frases devem fazer sentido mesmo que as express�es
de cita��o sejam removidas. Para exemplificar, n�o se deve usar:
\begin{quotation}
\dots conforme demonstrado por \cite{art-solimaes03}.
\end{quotation}
e sim:
\begin{quotation}
\dots conforme demonstrado por \citeasnoun{art-solimaes03}.

\dots conforme demonstrado na literatura
\citeaffixed{art-solimaes03}{e.g. }.\footnote{A abreviatura e.g.
significa por exemplo. Vem do latim \emph{exempli gratia}. Tamb�m se
usa, para o mesmo caso, v.g. (\emph{verbi gratia}) ou simplesmente
p.ex.}
\end{quotation}

\section{Considera��es finais}

Os simples fatos de utilizar corretamente uma boa ferramenta de
formata��o de textos e de seguir as recomenda��es quanto ao estilo
n�o garantem a qualidade do documento produzido. O principal
aspecto a ser levado em conta � a qualidade da reda��o.






%%
%% Cap�tulo 3: Express�es matem�ticas
%%

\mychapter{Express�es matem�ticas}
\label{Cap:matematica}

\LaTeX � insuper�vel no processamento de express�es matem�ticas. Express�es
simples como $2^{n}$ podem ser editadas no pr�prio texto. Equa��es
complexas como:
% eqnarray � o comando LaTeX de base para express�es multilinhas com
% alinhamento. O & indica o ponto onde todas as linhas devem ser
% alinhadas. Neste cap�tulo ser�o apresentadas outras op��es de alinhamento
\begin{eqnarray} \label{eq:PDF:RSR}
  p \left( \gamma \right) & = & \frac{1}{2} \sqrt{\frac{M}{\gamma \bar{\gamma}_{b}}} \frac{1}{ \prod_{i=1}^M {\sqrt{\tilde{\gamma}_i}}}
  \int_0^{\sqrt{M \delta}} \int_0^{\sqrt{M \delta} - r_M } \cdots
  \int_0^{\sqrt{M \delta} - \sum_{i = 3}^M {r_i } } \nonumber \\
  & & p \left( {\frac{\sqrt{M \delta} - \sum_{i = 2}^M {r_i }}{\sqrt{\tilde{\gamma}_1}} ,
  \frac{r_2}{\sqrt{\tilde{\gamma}_2}} , \ldots ,\frac{r_M}{\sqrt{\tilde{\gamma}_M}} } \right)
  \, dr_2 \cdots dr_{M-1} \, dr_M
\end{eqnarray}
% sem linha em branco
ou:
% sem linha em branco
\begin{equation} \label{eq:TrCGI}
  T(r) = \frac{1}{f_m}
  \left( \frac{\pi}{2} \sum_{i=1}^M
  {\tilde{r}_i^2 \dot{\varsigma}_i^2}\right)^{-1/2}
  \frac
  {\begin{array}{ll}
  \int_0^{\rho \sqrt{M}} \int_0^{\rho \sqrt{M} - r_M } \cdots
  \int_0^{\rho \sqrt{M} - \sum_{i = 3}^M {r_i } } \int_0^{\rho \sqrt{M} -
  \sum_{i = 2}^M {r_i } }  \\
  p \left( {\frac{r_1}{\tilde{r}_1} ,
  \frac{r_2}{\tilde{r}_2} , \ldots ,\frac{r_M}{\tilde{r}_M} } \right)
  \, dr_1 \, dr_2 \cdots dr_{M-1} \, dr_M \\ \end{array}}
  {\begin{array}{ll}
  \int_0^{\rho \sqrt{M}} \int_0^{\rho \sqrt{M} - r_M } \cdots
  \int_0^{\rho \sqrt{M} - \sum_{i = 3}^M {r_i } } \\
  p \left( {\frac{\rho \sqrt{M} - \sum_{i = 2}^M {r_i }}{\tilde{r}_1} ,
  \frac{r_2}{\tilde{r}_2} , \ldots ,\frac{r_M}{\tilde{r}_M} } \right)
  \, dr_2 \cdots dr_{M-1} \, dr_M \\ \end{array}}
\end{equation}

Pequenas express�es matem�ticas como $x_0^2$ podem ser inseridas
diretamente no texto, delimitadas por cifr�es (\texttt{\$}). Deve-se
evitar este recurso com express�es muito grandes, como

$\left[\begin{array}{cc} 1 & \frac{2}{x+1} \\ -2 &
1\end{array}\right]^{-1}$, porque o espa�amento entre as linhas fica
prejudicado. Para incluir express�es n�o numeradas maiores, pode-se
utilizar o par de delimitadores \verb|\[| e \verb|\]|, o que gera
express�es centralizadas na p�gina:

Um erro comum em express�es matem�ticas � o de digitar os nomes de
fun��es diretamente, sem utilizar os comandos apropriados. Por
exemplo, a express�o $\sin(\omega t+\phi)$ est� correta, enquanto que
a express�o $sin(\omega t+\phi)$ � interpretada pelo \LaTeX como sendo
o produto das vari�veis $s$, $i$ e $n$ pela express�o entre
par�nteses. 






%%
%% Cap�tulo 4: Figuras, gr�ficos e tabelas
%%

\mychapter{Figuras, tabelas e gr�ficos}
\label{Cap:figuras}

Uma das maiores dificuldades na edi��o de textos de qualidade � o
posicionamento dos elementos gr�ficos: figuras, gr�ficos e
tabelas. Como estes elementos muitas vezes s�o grandes, aparece o
dilema sobre o que fazer quando uma quebra de p�gina deveria acontecer
no meio do elemento. H� duas possibilidades:
\begin{enumerate}
\item O autor informa exatamente onde o elemento gr�fico deve ficar no
texto, evitando que quebras de p�ginas aconte�am no meio de um
elemento. O problema com esta abordagem � que todo o trabalho de
posicionamento pode ser perdido caso se inclua ou se exclua algum
texto ou elemento.
\item O editor de texto posiciona os elementos gr�ficos de forma a n�o
deixar espa�os em branco nas p�ginas. Estes elementos que podem ser
posicionados pelo editor s�o conhecidos como \emph{elementos
flutuantes}. O problema com esta abordagem � que o posicionamento
adotado pode n�o corresponder �s expectativas do autor.
\end{enumerate}

O \LaTeX\ oferece as duas possibilidades de posicionamento. Este
cap�tulo apresenta exemplos de inclus�o de elementos gr�ficos no
texto, bem como algumas ferramentas externas ao \LaTeX\ que podem ser
utilizadas para ger�-los.

\begin{figure}[tbp]
\caption{Trecho de \emph{Os Lus�adas}, de Luis de Cam�es}
\label{Fig:textoflutuante}
% hrule - linha horizontal
\hrule
% As minipage's s�o muito �teis para se colocar duas coisas na mesma linha
\begin{minipage}{0.45\linewidth}
% flushleft - alinha � esquerda
\begin{flushleft}
As armas e os bar�es assinalados\\
Que da ocidental praia lusitana\\
Por mares nunca dantes navegados\\
Passaram ainda al�m da Trapobana\\
Em perigos e guerras esfor�ados\\
Mais do que prometia a for�a humana\\
Entre gente remota edificaram\\
Novo reino, que tanto sublimaram
\end{flushleft}
\end{minipage}
\hfill
\begin{minipage}{0.45\linewidth}
% flushright - alinha � direita
\begin{flushright}
E tamb�m as mem�rias gloriosas\\
Daqueles reis que foram dilatando\\
A F�, o Imp�rio, as terras viciosas\\
De �frica e �sia andaram devastando,\\
E aqueles que por obras valorosas\\
Se v�o da lei da morte libertando:\\
Cantando espalharei por toda parte,\\
Se a tanto me ajudar o engenho e arte.
\end{flushright}
\end{minipage}
\hrule
\end{figure}

\begin{table}[bp]
% As minipage's s�o muito �teis para se colocar duas coisas na mesma linha
\begin{minipage}[b]{0.45\linewidth}
\begin{center}
\[
ax^2 + bx + c = 0
\]
\end{center}
\caption{Equa��o de segundo grau}
\label{Tab:equacaoflutuante}
\end{minipage}
\hfill
\begin{minipage}[b]{0.50\linewidth}
\begin{center}
\[
x = \frac{-b\pm\sqrt{b^2-4ac}}{2a}
\]
\end{center}
\caption{Ra�zes da equa��o da tabela~\ref{Tab:equacaoflutuante}}
\label{Tab:equacaoflutuante2}
\end{minipage}
\end{table}

� importante ressaltar que o que � numerado � o \texttt{caption} e n�o
a \texttt{figure} ou a \texttt{table}. Portanto, o \texttt{label} deve
ser colocado sempre ap�s o \texttt{caption} ao qual ele se
refere. Conforme ilustram as tabelas \ref{Tab:equacaoflutuante} e
\ref{Tab:equacaoflutuante2}, uma mesma \texttt{figure} ou
\texttt{table} pode ter mais de um ou nenhum \texttt{caption}.
O \texttt{caption} pode ser colocado antes do conte�do flutuante, como
na figura \ref{Fig:textoflutuante}, ou depois, como nas tabelas
\ref{Tab:equacaoflutuante} e \ref{Tab:equacaoflutuante2}. 

Em cada \verb|\begin{figure}| ou \verb|\begin{table}| pode-se incluir
um par�metro opcional com as op��es de posicionamento para este
elemento flutuante. Par�metros adicionais de comandos \LaTeX\ s�o
sempre fornecidos entre colchetes \texttt{[]}, enquanto os par�metros
obrigat�rios aparecem entre chaves \verb|{}|. As op��es dispon�veis
incluem as seguintes:
\begin{itemize}
\item[\tt h] O elemento pode ser posicionado na mesma posi��o em que ele
aparece no c�digo fonte do texto.
\item[\tt t] O elemento pode ser posicionado no topo de uma p�gina.
\item[\tt b] O elemento pode ser posicionado no fim de uma p�gina.
\item[\tt p] O elemento pode ser inclu�do em uma p�gina formada s� por
flutuantes.
\end{itemize}

\section{Tabelas em \LaTeX}
\label{Sec:tabelas}

Tabelas s�o constru�das com comandos pr�prios do \LaTeX, notadamente o
ambiente \texttt{tabular}. Nada obriga a que o ambiente
\texttt{tabular} esteja sempre posicionado em um elemento
flutuante. Se voc� quiser impor que uma tabela fique obrigatoriamente
em uma determinada posi��o do texto, basta n�o colocar o
\texttt{tabular} dentro de um \texttt{table}. Tabelas podem at� ser
inclu�das no meio de uma frase.  Por exemplo, eu posso dizer que se um
jogo da velha est� na configura��o \textsf{\tiny\begin{tabular}{c|c|c}
x & & x \\ \hline & & o \\ \hline x & o & \end{tabular}} e se o
jogador ``\textsf{x}'' sabe jogar, ent�o o jogador ``\textsf{o}'' ir�
perder, independentemente da jogada que fa�a.

 Exemplos de
colunas com diferentes larguras e alinhamentos podem ser vistos na
tabela \ref{Tab:larguracolunas}.

\begin{table}[htbp]
\begin{tabularx}{\linewidth}{|p{3cm}|X|l|} \hline
COLUNA p & COLUNA X & COLUNA l \\ \hline
Largura fixa (n�o depende do conte�do) &
Expand�vel &
Ajust�vel \\ \hline
Alinhada no topo &
Alinhada � esquerda &
Alinhada � esquerda \\ \hline
\end{tabularx}
\\[0.5cm]
\begin{tabularx}{\linewidth}{|b{3cm}|C|r|} \hline
COLUNA b & COLUNA C (ver \texttt{comandos.tex}) & COLUNA r \\ \hline
Largura fixa (n�o depende do conte�do) &
Expand�vel &
Ajust�vel \\ \hline
Alinhada na base &
Centralizada &
Alinhada � direita \\ \hline
\end{tabularx} 
\caption{Tabelas com colunas de diferentes larguras e alinhamentos}
\label{Tab:larguracolunas}
\end{table}

\section{Figuras em \LaTeX}
\label{Sec:figuras}

As figuras (imagens, desenhos, gr�ficos, etc.) devem ser produzidas
por ferramentas externas ao \LaTeX, salvas em um arquivo e inseridas
no texto usando o comando \texttt{includegraphics}. Da mesma forma
que as tabelas, as figuras podem ser flutuantes, caso sejam
inseridas dentro de um ambiente \texttt{figure}, ou ter uma posi��o
fixa no texto (como aqui: \includegraphics{figuras/eu}).

O formato em que voc� deve salvar os arquivos das figuras para que
possa inclu�-las no texto depende de como voc� pretende compilar
o c�digo fonte:
\begin{itemize}
\item se o texto vai ser compilado com \texttt{latex}, todos os
arquivos devem estar no formato EPS (\emph{Encapsulated PostScipt});
\item se o texto vai ser compilado com \texttt{pdflatex}, os
arquivos devem estar nos formatos PDF ou JPEG (outros formatos s�o
aceitos, mas estes s�o os recomend�veis).
\end{itemize}

A figura \ref{Fig:belmonte} mostra um exemplo de inclus�o de uma
imagem no texto \LaTeX.

\begin{figure}[htbp!] \begin{center}
% fbox faz uma borda ao redor do seu argumento
\fbox{\includegraphics[width=0.75\linewidth]{figuras/belmonte}}
\caption{Exemplo de imagem real}
\label{Fig:belmonte}
\end{center} \end{figure}

 

% Conclus�es
%%
%% Cap�tulo 5: Conclus�es
%%

\mychapter{Conclus�o}
\label{Cap:conclusao}

Neste cap�tulo fa�a o fechamento do seu trabalho. Informe como foram os resultados obtidos. Quais seus coment�rios e conhecimento obtidos a partir das experi�ncias executadas. Evite citar termos com ju�zo de valor como: "� ruim ...", "� bom ...", etc. prefira colocar: "foi poss�vel observar que 90\% das amostras ...". Aqui tamb�m n�o se deve colocar figuras, tabelas, etc. mas referenciar as que foram abordadas no texto fazendo o fechamento dos resultados.

%TRAB. FUTUROS
\section{Trabalhos Futuros}
\label{Sec:Trabalhos Futuros}
Descreva os trabalhos que voc� percebeu que daria pra fazer como melhoria, continuidade do que voc� fez e que por algum motivo n�o foi poss�vel de faz�-lo.

 

\backmatter


% Refer�ncias bibliog�ficas (geradas automaticamente)
%\bibliographystyle{ppgee}
\bibliographystyle{dcu}
\addcontentsline{toc}{chapter}{Refer�ncias bibliogr�ficas}
%Arquivo com as bibliografias
\bibliography{bibliografia}


\end{document}
